\section{第十轮拷打：优化记录总表与实施证据}

\begin{summaryBox}{这一章为什么必须单独写}
这一章不是为了显得内容多，而是为了把“设计想法、代码落地、测试证据、面试话术、未完成边界”放在同一个地方。面试里最怕的不是不会讲概念，而是把设计稿讲成生产能力。这里的表格就是为了防止这个问题：每一项都必须能回答“现在到底做到哪一步，证据是什么，没做到的边界在哪里”。
\end{summaryBox}

\begin{riskBox}{真实性原则}
本章的状态只有三类：\term{已落地}、\term{设计已定}、\term{未完成}。\term{已落地}必须有代码或测试证据；\term{设计已定}表示方案已经明确但代码还没完整接入；\term{未完成}表示面试时只能讲路线图，不能讲成当前能力。
\end{riskBox}

\begin{longtable}{>{\bfseries}p{0.16\textwidth}p{0.18\textwidth}p{0.16\textwidth}p{0.25\textwidth}p{0.17\textwidth}}
\toprule
主题 & 当前状态 & 证据 & 面试表达 & 边界 \\
\midrule
\endfirsthead
\toprule
主题 & 当前状态 & 证据 & 面试表达 & 边界 \\
\midrule
\endhead
Harness v2 状态 & 已落地 & \term{app/governance/harness.py}，\term{tests/governance/test\_harness.py} & Harness 不再只是最小快照器，已经有 v2 state、session config、backup、lease 和 summary。 & 仍不是完整分布式调度平台。 \\
Harness 原子写入与状态锁 & 已落地 & \term{harness-tasks.json.bak}、\term{.harness-state.lock} 测试 & 长任务状态文件写入必须避免半写入；写前备份、临时文件、原子替换和轻量 transaction lock 是最低治理线。 & 仍不是完整多 worker scheduler。 \\
Harness JSON 恢复 & 已落地 & corrupt JSON backup recovery 测试 & state 文件是恢复上下文，损坏时必须有明确恢复路径，不能静默失败。 & backup 也损坏时会停止，不能自动重建任务语义。 \\
Harness active marker & 已落地 & marker 创建和 idle 清理测试 & active marker 用来区分 harness 是否接管任务生命周期，避免无任务时 hook 误触发。 & marker 不是租约锁，不能替代 worker coordination。 \\
Harness lease & 已落地 & expired lease 测试 & lease 解决“谁正在做”的事实问题；没有 lease，恢复时无法区分正常执行和中断。 & 还没有多 worker 抢租约和工作目录隔离。 \\
Harness summary 与依赖校验 & 已落地 & summary、cycle、blocked propagation 测试 & supervisor 必须能给出任务总览，并在循环依赖或失败依赖阻塞时显式标记 terminal failed。 & 依赖校验是安全网，不是完整任务调度器。 \\
MCP allowlist & 已落地 & \term{McpPolicyProxy} 与 untrusted/tool 测试 & MCP 只解决协议，不解决信任；server 和 tool 必须先 allowlist。 & Planner 还未默认大量生成 MCP 节点。 \\
MCP schema 校验 & 已落地 & nested schema 与 array item 测试 & tool inputSchema 是执行前最后一道结构化约束，不能只靠模型理解参数；当前已覆盖 object、array、enum、required 和 additionalProperties 子集。 & 仍不是完整 JSON Schema 引擎。 \\
MCP 写操作审批 & 已落地 & write tool approval 与 raw args hash 测试 & 非只读或写工具必须返回 approval request；approval id 要绑定真实参数指纹，避免不同写入意图复用旧审批。 & 审批最终裁决仍要落到本地 DB 聚合根。 \\
MCP redaction 与稳定审计 hash & 已落地 & configured redaction keys、stable result hash 测试 & MCP 返回值进入上下文前必须递归脱敏，并记录稳定 result hash 与 fingerprint；hash 不应受字典字段顺序影响。 & redaction key 需要随业务补充。 \\
Skill scope 保护 & 已落地 & scope metadata binding 与 manual enable scope 测试 & 手动启用不能绕过 tenant/project scope；metadata 也必须强约束 global/tenant/project 字段绑定。 & 还没有独立的企业级 RBAC UI。 \\
Skill top-k & 已落地 & top-k truncation 测试 & Skill 不是匹配越多越好，最终进入 prompt 的能力必须有上限。 & top-k 仍是规则配置，不是学习排序。 \\
Skill snapshot & 已落地 & resolved snapshot version 测试 & session 必须绑定 skill version 和 prompt snapshot，避免运行中行为漂移。 & 还没有版本回收生命周期管理 UI。 \\
Agent tracing & 设计已定 & 现有 agent trace 与 tool histories & 面试时讲 tracing 是为了解释“为什么走到这里”，不是为了展示日志很多。 & 仍需更完整的 trace correlation。 \\
Agent eval & 设计已定 & metrics 目录与 RAG/browser collectors & Agent eval 要同时看任务成功、证据质量、工具成本和恢复能力。 & 尚未形成统一 eval dashboard。 \\
多 worker Harness & 未完成 & 无完整调度测试 & 多 worker 需要 state lock、lease、独立 worktree、rollback 隔离，不是简单开多个进程。 & 当前只能讲路线图。 \\
Skill 自定义执行器 & 未完成 & 无 DAG node template 实现 & Skill 当前影响 prompt 和工具边界；更高阶是自定义节点模板与执行策略。 & 面试不能说已经具备。 \\
MCP 主流量化 & 未完成 & Planner 仍不大量生成 MCP 节点 & 当前 MCP 是安全接入通路；主业务默认流量仍以 search/browser/rag 为主。 & 需要更多 server registry 和 planner policy。 \\
\bottomrule
\end{longtable}

\section{第十一轮拷打：Harness 每一处细节怎么拆}

\begin{summaryBox}{一句话定位 Harness}
Harness 是长任务 Agent 的任务级控制面。它不替代 LangGraph 的业务编排，也不替代数据库的审批、预算和审计事实源；它负责把跨上下文、跨进程、可恢复、可重试、可诊断的任务状态从短生命周期请求里抽出来。
\end{summaryBox}

\subsection{为什么 Harness 不是“再加一个日志文件”}
\begin{analysisBox}{核心判断}
如果一个 Agent 任务能在一次请求、一次上下文窗口、一次进程生命周期里完成，那么普通日志和内存状态可能够用。但深度研究任务不是这样：它会搜索、浏览、RAG、分析、反思、重规划，还可能等待审批、遇到预算熔断、发生浏览器超时、跨上下文恢复。此时日志不是附属物，任务状态也不是 UI 展示字段；它们是下一次执行能不能继续的事实依据。

Harness 的价值就在这里：它把“任务现在在哪里、谁在执行、执行到哪个 checkpoint、失败是否可重试、等待哪个审批、预算用了多少、最近一次错误是什么”这些信息外置成可读取的控制面。这样当 agent 上下文被压缩、服务重启、worker 中断时，系统不是靠口头回忆恢复，而是读取 state file、progress log 和 DB runtime state 重新判断。
\end{analysisBox}

\subsection{Progress file 为什么是 append-only}
\begin{qaBox}{面试回答}
progress file 必须 append-only，因为它承担的是审计和恢复线索，而不是当前状态视图。当前状态可以覆盖写，但历史动作不能覆盖写；否则一旦恢复失败，你无法知道是哪个 worker、哪个 checkpoint、哪个错误分类导致了当前状态。
\end{qaBox}

\begin{analysisBox}{深入分析}
append-only 的价值不是“日志看起来完整”，而是能建立事件序列。长任务系统里很多 bug 不是静态状态能解释的，比如任务从 running 变成 retryable failed，再变成 running，最后又 awaiting approval。只看最终 JSON，你只能知道它现在暂停了；看 progress file，你才能知道中间经历了几次重试、是否发生 rollback、是否重复进入同一个失败分支。

这也是为什么日志格式要 grep-friendly。错误类型、session id、task id、checkpoint 和 category 要在同一行，而不是散在多行自然语言里。面试官如果追问“为什么标准化日志格式”，答案不是方便人看，而是方便 agent 和脚本恢复时快速定位。
\end{analysisBox}

\subsection{Tasks JSON 为什么是 structured state}
\begin{qaBox}{面试回答}
progress file 负责事件历史，tasks JSON 负责当前可执行状态。恢复时真正驱动决策的是 tasks JSON：哪些任务 pending，哪些 in progress，哪些 failed 但可重试，哪些 blocked，哪个 worker 持有 lease，checkpoint 在哪里。
\end{qaBox}

\begin{analysisBox}{字段拆解}
最小可用的 tasks JSON 至少要有六类字段。第一类是身份字段，例如 task id、title、priority。第二类是生命周期字段，例如 status、runtime status、attempts、max attempts、completed at。第三类是依赖字段，例如 depends on，用来判断 blocked propagation。第四类是恢复字段，例如 started at commit、checkpoint、state snapshot。第五类是 worker 字段，例如 claimed by、lease expires at、last heartbeat。第六类是验证字段，例如 validation command、timeout 和 on failure cleanup。

DeepIntel 当前的 Harness 更偏研究 session supervisor，所以它把 checkpoint、budget、pending approval 和 state snapshot 作为主要状态；这和通用 harness 技能里的 validation command 语义不同。面试时要讲清楚：项目里的 Harness 是研究任务控制面，不是通用 CI runner。
\end{analysisBox}

\subsection{Active marker 的意义}
\begin{qaBox}{面试回答}
active marker 是开关，不是锁。它的作用是告诉 hook 或外部 supervisor：“当前目录有 harness 任务需要接管”。没有 marker 时相关 hook 应该 no-op，避免普通项目操作被 harness 误拦截。
\end{qaBox}

\begin{analysisBox}{容易混淆的点}
active marker 不能替代 lease，也不能替代 lock。marker 回答“有没有 harness 任务处于活跃状态”，lease 回答“这个任务现在由谁执行”，lock 回答“谁能修改共享状态文件”。这三个问题不同，不能用一个文件同时解决。面试时把这三个概念拆开，会比简单说“有个 .harness-active 文件”可信得多。
\end{analysisBox}

\subsection{Lease 为什么是长任务 Agent 的核心字段}
\begin{qaBox}{面试回答}
lease 解决的是所有权和超时恢复。没有 lease，新 session 看到一个 in progress 任务时无法判断它是仍在运行，还是上一个 worker 已经死掉。lease 过期后才有资格进入恢复逻辑，否则并发 worker 会互相踩状态。
\end{qaBox}

\begin{analysisBox}{租约设计的三个细节}
第一，lease 必须有稳定 worker id，不能只靠进程号。进程号在容器、重启和不同机器上都不稳定，worker id 至少要能解释“哪个执行者认领了任务”。第二，lease 必须续约，而不是只在开始时写一次。长任务执行过程中每个 checkpoint 都应更新 heartbeat，否则一个正常运行但耗时很长的任务会被误判为 stale。第三，lease 过期只代表可以恢复，不代表可以无脑覆盖。恢复前仍要检查 checkpoint、DB runtime state、pending approval 和最近错误。
\end{analysisBox}

\subsection{为什么恢复边界不能只靠 LangGraph checkpointer}
\begin{qaBox}{面试回答}
LangGraph checkpointer 保存的是图内状态快照；Harness 负责的是任务级控制决策。checkpointer 能告诉你图走到哪里，不能单独告诉你是否还有预算、审批是否通过、worker 租约是否过期、是否应该重试或终止。
\end{qaBox}

\begin{analysisBox}{职责边界}
把恢复问题全部交给 checkpointer，会漏掉任务级语义。比如一个节点已经写入 checkpoint，但它等待审批；如果恢复时只从 checkpoint 继续跑，而不查本地审批聚合，就可能绕过人工确认。再比如浏览器节点失败三次，checkpoint 里可能记录了失败结果，但是否还有 retry budget、是否需要 rollback、是否应该进入 terminal failed，是 supervisor 的判断。

所以更稳的架构是：checkpointer 负责“状态在哪里”，Harness 负责“能不能继续、从哪里继续、由谁继续”，DB 负责“审批、预算、审计的事实源”。三者叠加才是完整恢复，不是某一个组件包打天下。
\end{analysisBox}

\subsection{Rollback 为什么不能随便做}
\begin{riskBox}{高风险点}
通用 harness 技能里会提到失败后 \term{git reset --hard} 和 \term{git clean -fd}。这是对隔离工作树的任务 runner 才安全。DeepIntel 当前是在共享工作区里协作，不能把这种 rollback 当作默认行为，否则可能删除用户改动。
\end{riskBox}

\begin{analysisBox}{面试表达}
正确表达是：rollback 需要满足两个条件。第一，任务必须运行在独立工作树或独立 clone 里，避免清理其他 worker 或用户改动。第二，任务开始时必须记录 base commit，失败后验证这个 commit 仍存在。没有这两个条件，自动 rollback 就不是安全恢复，而是破坏性操作。

这也是为什么当前项目里更适合把 Harness 定位成 session supervisor，而不是完全自动化代码任务 runner。它能记录 checkpoint、lease、预算和审批状态，但不会在共享工作区里随意 reset。
\end{analysisBox}

\subsection{Dependency 与 blocked propagation}
\begin{qaBox}{面试回答}
依赖传播解决的是“下游任务是否还有意义”。如果一个 P0 任务永久失败，依赖它的任务应该被标记 blocked 或 failed with dependency reason，而不是一直 pending 让 worker 反复挑选。
\end{qaBox}

\begin{analysisBox}{工程意义}
没有 blocked propagation 的任务队列会出现假进度：pending 数量看起来还有很多，但其中大部分永远不能执行。成熟的 supervisor 要能区分“还没轮到的 pending”和“因为上游失败永远不能跑的 blocked”。这个区别不仅影响调度，也影响面试里对系统可靠性的解释。
\end{analysisBox}

\subsection{Stats 为什么是恢复协议的一部分}
\begin{qaBox}{面试回答}
stats 不只是 UI 汇总，它是每个 session 结束时对任务状态的一次结构化结算。它让下一次恢复能快速知道 completed、failed、pending、blocked、attempts、checkpoint 的整体分布。
\end{qaBox}

\begin{analysisBox}{为什么要写进 progress}
stats 写进 progress file 的好处是可审计。即使 tasks JSON 后续被更新，你仍能看到每个 session 结束时的全局状态。这对于长任务尤其重要，因为问题常常不是某一个 task 错了，而是“系统什么时候开始进入大量 retryable failed、什么时候出现 blocked、什么时候 session count 打满”。
\end{analysisBox}

\subsection{Harness 面试追问清单}
\begin{longtable}{>{\bfseries}p{0.22\textwidth}p{0.72\textwidth}}
\toprule
追问 & 稳妥回答 \\
\midrule
\endfirsthead
\toprule
追问 & 稳妥回答 \\
\midrule
\endhead
为什么不把 Harness 放进 LangGraph 图里？ & 因为图表达业务执行，Harness 表达任务级控制。把 lease、审批、预算、恢复都塞进图，会让图变成混杂控制器。 \\
progress 和 tasks 为什么要分两个文件？ & 一个是事件历史，一个是当前状态。历史要 append-only，当前状态要可覆盖、可恢复、可查询。 \\
lease 过期后能直接接管吗？ & 不能直接覆盖。过期只代表有资格恢复，还要检查 checkpoint、DB runtime state、pending approval 和最近错误。 \\
为什么 active marker 不是锁？ & marker 只表达 harness 是否活跃，不表达谁持有状态写权限，也不表达任务所有权。 \\
多 worker 最难在哪里？ & 难在共享状态锁、任务租约、独立工作树、rollback 隔离和幂等提交，不是多开几个进程。 \\
checkpoint 和 snapshot 有什么区别？ & checkpoint 是恢复点，snapshot 是当时的完整状态材料。checkpoint 更偏指针，snapshot 更偏事实内容。 \\
为什么 premature completion 是大风险？ & Agent 容易把“完成了一部分步骤”误报成“目标已完成”。结构化 task criteria 和 validation 是防线。 \\
验证命令缺失怎么办？ & 不能宣布 completed。应该记录 CONFIG 错误并停止，避免没有客观校验的完成。 \\
JSON 损坏怎么办？ & 先恢复 .bak；没有有效备份就停止。不能从日志凭空重建 validation、depends on 和 cleanup。 \\
为什么日志要 grep-friendly？ & 因为恢复和排障时需要脚本快速筛 ERROR、SESSION、CHECKPOINT、RECOVERY，而不是靠人读散文。 \\
\bottomrule
\end{longtable}

\section{第十二轮拷打：MCP 协议与治理怎么讲深}

\begin{summaryBox}{一句话定位 MCP}
MCP 是外部能力接入协议，不是安全治理系统。它让 client 和 server 用统一方式描述资源、提示和工具，但“该不该信任这个 server、该不该允许这个 tool、参数是否合法、返回值是否脱敏、写操作是否需要审批”仍然是应用自己的责任。
\end{summaryBox}

\subsection{MCP 的三层结构}
\begin{analysisBox}{协议层、传输层、治理层}
第一层是消息层，核心是 JSON-RPC 2.0。无论本地 stdio 还是远程 Streamable HTTP，本质都是 client 发 method、params、id，server 回 result 或 error。第二层是传输层，本地常见 stdio，远程推荐 Streamable HTTP。第三层是治理层，这一层不是 MCP 自动提供的，而是 DeepIntel 通过 \term{McpPolicyProxy} 补上的。

面试时不要把这三层混在一起。说“我们用了 MCP，所以安全”是不严谨的；更准确的说法是“我们按 MCP 接入外部能力，但所有调用先经过 policy proxy 做 server trust、tool allowlist、schema、审批和 redaction”。
\end{analysisBox}

\subsection{为什么 JSON-RPC 适合 MCP}
\begin{qaBox}{面试回答}
JSON-RPC 适合 MCP，是因为它轻量、跨语言、可读、调试成本低。MCP server 可能用 Python、TypeScript、Go 或其他语言写，JSON-RPC 给它们一个共同的调用结构。
\end{qaBox}

\begin{analysisBox}{不是因为先进，而是因为合适}
协议设计的价值不总是越复杂越好。MCP 的核心诉求是让模型应用发现和调用工具，而不是构建一个重型微服务网关。JSON-RPC 的 method、params、result、error 已经足够表达 tools/list、tools/call、resources/read 这类交互。它不像 REST 那样把语义拆到 URL 和 HTTP verb，也不像 gRPC 那样需要 IDL 和二进制协议，对 agent 工具生态来说门槛更低。
\end{analysisBox}

\subsection{stdio 与 Streamable HTTP 怎么比较}
\begin{longtable}{>{\bfseries}p{0.18\textwidth}p{0.24\textwidth}p{0.24\textwidth}p{0.24\textwidth}}
\toprule
维度 & stdio & Streamable HTTP & 面试落点 \\
\midrule
\endfirsthead
\toprule
维度 & stdio & Streamable HTTP & 面试落点 \\
\midrule
\endhead
部署位置 & 本地子进程 & 远程服务 & 本地工具优先 stdio，企业共享能力更适合 HTTP。 \\
暴露面 & 无网络端口 & 有网络入口 & HTTP 必须认真做 auth、trust 和 rate limit。 \\
调试 & 看 stdin/stdout & 看 HTTP 请求与 SSE 流 & 都是 JSON-RPC，消息层不变。 \\
扩缩容 & 依赖本地进程 & 可放到服务侧扩展 & 远程 server 更适合团队共享。 \\
安全风险 & 本地文件/进程权限 & 网络、身份、租户、跨域 & 不同传输有不同治理重点。 \\
\bottomrule
\end{longtable}

\subsection{Tool inputSchema 为什么必须落到执行前}
\begin{qaBox}{面试回答}
inputSchema 是 Agent 工具调用从自然语言走向真实执行的第一道确定性边界。模型可以理解错参数，但 schema 不能跟着理解错；参数不符合 type、required、enum、additionalProperties，就不应该执行。
\end{qaBox}

\begin{analysisBox}{DeepIntel 的实现边界}
当前项目实现的是轻量 JSON Schema 子集，覆盖 object、properties、required、additionalProperties、基础类型和 enum。这个选择是有意的：它足以挡住大多数错误参数和意外字段，又不会引入完整 schema 引擎的复杂度。面试时可以说“当前是工程上足够用的 schema 子集，不是完整 JSON Schema 解释器”。
\end{analysisBox}

\subsection{为什么 tool annotations 不能当安全依据}
\begin{riskBox}{关键风险}
MCP server 可能给 tool 标注 read-only、destructive、open-world 等信息，但这些 annotation 只能作为提示，不能作为最终安全裁决。因为 server 本身就是外部信任边界的一部分，不能让它自己声明自己安全，然后应用就无条件相信。
\end{riskBox}

\begin{analysisBox}{正确做法}
正确做法是应用侧维护 registry：server 是否 trusted、允许哪些 tools、哪些 tools 是 write tools、是否 read-only default、需要哪些 approval。server 提供的 schema 和 annotation 可以进入 registry 审核流程，但最终执行时要以本地 registry 为准。这和审批系统的 source of truth 原则是一致的：外部提供信号，本地做裁决。
\end{analysisBox}

\subsection{MCP authorization 怎么回答}
\begin{qaBox}{面试回答}
远程 MCP 不能只靠“知道 endpoint”来调用，必须有 authorization。规范层面会涉及 OAuth 语义，但工程落地时还要做 server registry、token 管理、scope、租户隔离和审计。授权解决“你是谁”，policy proxy 解决“你能做什么、这次参数能不能执行”。
\end{qaBox}

\begin{analysisBox}{授权和治理的区别}
Authorization 不是全部治理。一个用户有权访问某个 MCP server，不代表他可以调用这个 server 的所有工具；有权调用某个工具，也不代表任意参数都合法；参数合法，也不代表返回值可以原样进入模型上下文。面试里要把这条链路拆开：身份认证、工具授权、参数校验、审批、执行、返回脱敏、审计。
\end{analysisBox}

\subsection{MCP 面试追问清单}
\begin{longtable}{>{\bfseries}p{0.22\textwidth}p{0.72\textwidth}}
\toprule
追问 & 稳妥回答 \\
\midrule
\endfirsthead
\toprule
追问 & 稳妥回答 \\
\midrule
\endhead
MCP 是不是等于工具治理？ & 不是。MCP 是协议，治理要在应用侧做 policy proxy。 \\
为什么不用直连 MCP tool？ & 直连会丢失 server trust、tool allowlist、schema、审批和 redaction。 \\
MCP 底层消息是什么？ & JSON-RPC 2.0，method、params、id、result/error。 \\
本地 MCP 为什么常用 stdio？ & 配置简单、延迟低、没有网络暴露面，适合本地工具。 \\
远程为什么用 Streamable HTTP？ & 单端点承接请求和响应，需要流式时返回 SSE，比双端点 SSE 更简单。 \\
inputSchema 怎么用？ & 执行前校验参数结构，不通过就返回 schema invalid。 \\
写工具怎么处理？ & 默认需要 approval request，本地审批聚合通过后才执行。 \\
返回值为什么要 redaction？ & MCP server 返回可能含 token、内部 ID、secret，进入上下文前必须脱敏。 \\
server fingerprint 有什么用？ & 审计时证明调用的是哪个 server 版本或实例，避免 tool list 漂移无法追责。 \\
Planner 为什么还没大量生成 MCP 节点？ & 当前 MCP 是安全接入通路，主流量仍以 search/browser/rag 为主，MCP 主流量化是后续演进。 \\
\bottomrule
\end{longtable}

\section{第十三轮拷打：Agent 开发岗位底层能力}

\begin{summaryBox}{岗位视角}
Agent 开发岗位不是只会写 prompt。真正会被追问的是运行时、状态、工具、权限、检索、评测、成本、观测和失败恢复。下面按平台基础设施、应用/RAG研究、安全治理三条线拆。
\end{summaryBox}

\subsection{平台基础设施线}
\begin{longtable}{>{\bfseries}p{0.24\textwidth}p{0.30\textwidth}p{0.38\textwidth}}
\toprule
底层问题 & 应该掌握什么 & DeepIntel 怎么讲 \\
\midrule
\endfirsthead
\toprule
底层问题 & 应该掌握什么 & DeepIntel 怎么讲 \\
\midrule
\endhead
运行时状态 & state machine、checkpoint、thread id、resume & ResearchState 是业务状态，Harness 是任务级状态，DB 是事实状态。 \\
并发调度 & DAG batch、fan-out、join、依赖 & Planner 产 DAG，LangGraph 按可执行批次分发 Search/Browser/RAG。 \\
长任务恢复 & checkpoint、lease、heartbeat、retry budget & Harness 记录 lease、checkpoint 和 progress，避免上下文丢失后无法恢复。 \\
工具执行 & tool schema、timeout、sandbox、error category & 工具调用先过确定性规则，再进入真实资源。 \\
HITL & interrupt、approval、resume、idempotency & 动作审批以 DB 为权威源，外部系统只是信号。 \\
观测性 & trace id、span、tool history、event stream & SSE 推进 UI，tool history 和 audit 解释每次调用。 \\
成本控制 & token budget、tool calls、wall clock、hard stop & session budget state 和 breaker 控制长任务膨胀。 \\
错误语义 & retryable vs terminal vs awaiting approval & runtime status 和 public status 分离，避免 UI 状态污染执行语义。 \\
存储设计 & session、checkpoint、audit、budget、skill meta/content & PostgreSQL 保持事实源，Redis/缓存只做加速。 \\
扩展边界 & plugin、skill、MCP、custom nodes & Skill 是领域能力层，MCP 是外部工具协议，二者不是同一层。 \\
\bottomrule
\end{longtable}

\subsection{应用/RAG 研究线}
\begin{longtable}{>{\bfseries}p{0.24\textwidth}p{0.30\textwidth}p{0.38\textwidth}}
\toprule
底层问题 & 应该掌握什么 & DeepIntel 怎么讲 \\
\midrule
\endfirsthead
\toprule
底层问题 & 应该掌握什么 & DeepIntel 怎么讲 \\
\midrule
\endhead
Query decomposition & 子问题、依赖、证据源选择 & Planner 不是写答案，而是生成可执行研究 DAG。 \\
Hybrid retrieval & BM25、vector、rerank、metadata filter & RAG 负责内部证据，不替代 Web 取证。 \\
Citation grounding & evidence id、claim attribution、coverage & Report 必须用 citation 回溯证据，Reflection 检查覆盖率。 \\
Browser extraction & accessibility、DOM、readability、fallback & Browser 节点不是简单抓 HTML，而是多策略提取。 \\
Research quality & faithfulness、answer relevance、coverage & 质量评估要看证据充分性，不只看文风。 \\
Replanning & failure hints、known bad paths、budget & Reflection 不通过时进入 replan，但受预算和失败记忆约束。 \\
Internal-first & 内部资料优先、web 补充 & DeepIntel 默认 internal-first，避免一上来就联网扩大成本和风险。 \\
Report synthesis & structure、引用、冲突证据、边界 & Analyst 综合，Report 输出，边界和不确定性要写明。 \\
Eval dataset & query set、expected evidence、metrics & metrics 目录已有实验基础，但统一 dashboard 未完成。 \\
成本/质量权衡 & top-k、depth、browser budget、model choice & 研究深度不是无限扩张，要有 budget profile。 \\
\bottomrule
\end{longtable}

\subsection{安全治理线}
\begin{longtable}{>{\bfseries}p{0.24\textwidth}p{0.30\textwidth}p{0.38\textwidth}}
\toprule
底层问题 & 应该掌握什么 & DeepIntel 怎么讲 \\
\midrule
\endfirsthead
\toprule
底层问题 & 应该掌握什么 & DeepIntel 怎么讲 \\
\midrule
\endhead
Prompt injection & instruction hierarchy、untrusted content、tool boundary & 浏览网页和 MCP 返回都属于不可信内容，不能让其改写系统策略。 \\
Excessive agency & least privilege、approval、read-only default & 高风险工具默认审批，MCP 非只读不直接执行。 \\
Data leakage & redaction、tenant scope、row limit & MCP 和 tool result 进入上下文前要脱敏，Skill scope 不能越权。 \\
Tool poisoning & server trust、tool allowlist、schema drift & MCP server 必须 registry 审核，fingerprint 进入审计。 \\
Supply chain & dependency、MCP server provenance、skill content & Skill 和 MCP 都是扩展面，需要版本和来源管理。 \\
Auditability & who/what/why/result hash & 工具调用必须记录调用者、参数范围、结果摘要、是否被拦。 \\
Sandbox & filesystem、network、CPU/memory、timeout & 真正有副作用的工具要在受限环境里跑。 \\
Approval SOT & DB aggregate、external signal、idempotency & 外部审批不能直接驱动执行，本地 DB 是最终裁决。 \\
Policy testing & unit tests、red-team cases、regression & 安全规则要用确定性测试证明，不靠 prompt 承诺。 \\
Incident response & pause、revoke、rollback、forensics & 当前有暂停/审批/审计基础，但全量法证平台未完成。 \\
\bottomrule
\end{longtable}

\subsection{Agent 岗高压题库：平台基础设施}
\begin{enumerate}[leftmargin=2em]
  \item 如果 LangGraph 已经有 checkpoint，为什么还要 Harness？回答时强调图内状态与任务级控制面不同。
  \item 任务恢复时先看 DB、Harness 还是 checkpointer？回答时讲三者职责：DB 是事实，Harness 是控制，checkpointer 是图状态。
  \item 多 worker 怎么避免抢同一个任务？回答 lease、state transaction lock、独立工作树。
  \item 为什么不能在共享仓库里自动 \term{git reset --hard}？回答 rollback 隔离和用户改动保护。
  \item retryable failed 和 terminal failed 怎么区分？回答错误类别、attempt budget、依赖传播。
  \item 为什么 public status 和 runtime status 要分开？回答 UI 友好状态不能污染执行语义。
  \item budget breaker 是软提示还是硬控制？回答必须是硬熔断加优雅降级。
  \item 工具超时算模型失败还是工具失败？回答是工具执行失败，进入 tool error category。
  \item Trace 和 audit 有什么区别？Trace 解释流程，audit 证明动作事实。
  \item Session snapshot 为什么要版本化？避免恢复时状态结构不兼容。
  \item 如果 checkpoint 存在但 approval 已过期怎么办？不能继续执行，必须重新进入审批或终止。
  \item 如何防止 Agent 伪造 tool result？Observation 必须来自工具层真实返回。
  \item 长任务怎么判断完成？不能靠模型说完成，要靠任务 criteria 和验证。
  \item SSE 断了任务是否停止？不应该，SSE 是交互通道，不是任务事实源。
  \item 如果 DB 写入成功但 Harness 写入失败怎么办？要记录错误并在恢复时以 DB 事实为准补偿。
  \item 为什么任务状态不能只放 Redis？Redis 适合缓存，不适合唯一事实源。
  \item 什么是 idempotency key？用于避免审批回调、工具执行或状态更新重复生效。
  \item 如何设计 stop/resume？stop 改 runtime state，resume 读取 checkpoint 和 harness snapshot。
  \item 如何做 model fallback？fallback 不能破坏结构化输出 schema 和工具调用契约。
  \item Agent 失败后如何生成部分报告？基于已有证据输出，并标注覆盖不足和未执行节点。
  \item 为什么 orchestration 不是 prompt engineering？因为它处理状态、依赖、恢复和副作用。
  \item 如何把任务拆成 DAG？节点是可执行取证动作，边表达依赖，不是自然语言段落。
  \item 什么情况下 replan？Reflection 发现证据不足、失败路径重复、原 DAG 不可达。
  \item 如何限制 replan 膨胀？max revision、budget、known failures 和 node cap。
  \item 如果 worker 死在工具执行中怎么办？lease 过期后恢复，检查副作用审计和 checkpoint。
\end{enumerate}

\subsection{Agent 岗高压题库：应用与 RAG}
\begin{enumerate}[leftmargin=2em]
  \item Deep Research 和普通 RAG 的核心差异是什么？回答任务状态机 vs 上下文增强回答。
  \item 为什么 internal-first？回答成本、安全和内部证据优先级。
  \item BM25 和向量检索怎么互补？BM25 抓关键词，向量抓语义。
  \item 为什么还要 rerank？召回集合和最终上下文选择不是同一个问题。
  \item citation 怎么保证可信？证据 id、引用覆盖、claim 对齐和 reflection。
  \item 网页内容怎么防 prompt injection？网页是证据，不是指令。
  \item Browser 提取失败如何降级？accessibility、readability、raw text、snippet fallback。
  \item 证据冲突怎么办？不强行合并，报告冲突来源和可信度。
  \item RAG 命中后还要不要 Web？看任务类型、证据时效性和用户允许策略。
  \item 如何评价研究报告？faithfulness、coverage、answer relevance、citation precision。
  \item 为什么不能只看 BLEU/ROUGE？研究任务更关注事实和证据，不是字面重合。
  \item 如何构建 eval dataset？按任务类型、证据源、难度、时间敏感性分层。
  \item 如何处理时间敏感问题？必须重新取证，不能靠模型记忆。
  \item 如何减少 hallucination？证据优先、拒答、引用、reflection、工具真实 observation。
  \item Planner 输出 JSON 失败怎么办？schema 校验、repair、fallback DAG。
  \item fallback DAG 会不会太死？它是最低可执行骨架，不是替代高级规划。
  \item 什么是 evidence coverage？关键 claim 是否都有证据支持。
  \item 如何避免检索上下文过长？top-k、rerank、摘要、chunk budget。
  \item 报告应该输出多长？由 budget profile 和任务复杂度控制。
  \item 为什么 Analyst 和 Report 分开？分析是证据综合，报告是表达和引用组织。
  \item Reflection 是否会增加成本？会，所以要受 revision budget 控制。
  \item 搜索结果不可靠怎么办？多源交叉、来源评级、Browser 深读。
  \item 内部资料过期怎么办？标注时间，必要时 Web 补证。
  \item 如何处理用户给的错误前提？先校正前提，再回答条件化结论。
  \item 什么时候不应该联网？用户禁止、内部证据足够、风险或预算不允许。
\end{enumerate}

\subsection{Agent 岗高压题库：安全治理}
\begin{enumerate}[leftmargin=2em]
  \item Prompt injection 的本质是什么？不可信内容试图越权改变上层指令。
  \item 为什么 tool schema 比 prompt 安全？schema 是确定性约束，prompt 是概率性引导。
  \item MCP tool 为什么默认不可信？server 是外部边界，tool list 可能漂移。
  \item read-only default 有什么意义？把副作用作为例外，而不是默认。
  \item 写操作审批如何防重复执行？approval id、action id、DB version、幂等键。
  \item 为什么外部审批只能是信号？最终执行发生在本地系统，本地 DB 必须裁决。
  \item 数据脱敏在哪做？进入模型上下文和审计前都要做。
  \item 如果 redaction 漏了怎么办？补充 policy key、回归测试、审计影响范围。
  \item 如何防止 skill 越权？scope 过滤、manual enable 不绕过、session snapshot。
  \item Skill 内容是不是也可能被投毒？是，所以要版本、来源、启用状态和审核。
  \item Agent 能不能直接写 SQL？不应该，应该走参数化模板和受控 DB 接口。
  \item 错误信息为什么要净化？原始错误可能泄露 schema、路径、secret 或内部拓扑。
  \item 审计日志要记录什么？调用者、tool、参数摘要、结果 hash、策略裁决、时间。
  \item 如何处理高风险工具？默认隐藏或只读，必要时审批和沙箱。
  \item 什么是 excessive agency？Agent 被授予超出任务需要的权限和自主性。
  \item 如何限制 excessive agency？least privilege、budget、approval、tool allowlist。
  \item MCP authorization 和 allowlist 区别？auth 证明身份，allowlist 限制能力。
  \item 如何做红队测试？注入、越权、敏感字段、schema 绕过、审批重放。
  \item 安全策略失败时返回什么？稳定 safe error，不暴露内部细节。
  \item 为什么不能相信模型自我审查？模型审查是辅助，执行面必须确定性控制。
  \item Tool result 能不能带指令？可以出现文本，但应作为数据处理，不能提升为指令。
  \item 多租户下最大风险是什么？数据越界和 skill/tool scope 漂移。
  \item 如何处理删除或更新的 skill？新 session 不再匹配，旧 session 继续使用版本 snapshot。
  \item 审批过期怎么办？执行前检查本地状态和 expiry，过期则拒绝或重新申请。
  \item 如何解释“安全不是阻碍能力”？治理让能力在可控边界内扩大，否则工具越多风险越大。
\end{enumerate}

\section{第十四轮拷打：Skill 体系深挖}

\begin{summaryBox}{一句话定位 Skill}
Skill 是领域能力的系统级扩展点，不是多拼一段 prompt。它要能被发现、匹配、启用、禁用、版本化、审计，并能影响 Planner、Agent hints 和工具 allowlist。
\end{summaryBox}

\subsection{为什么 Markdown + YAML frontmatter}
\begin{qaBox}{面试回答}
frontmatter 承载机器可校验字段，Markdown body 承载人类可编辑说明。这比纯 prompt 可治理，比纯 JSON/YAML 更适合人维护。
\end{qaBox}

\begin{analysisBox}{字段拆解}
frontmatter 至少要包含 name、slug、description、version、enabled、priority、trigger patterns、allowed tools、agent hints、tags、keywords、scope。body 里可以有 overview、prompt、constraints。这样 registry 能用 metadata 做匹配和治理，Agent 能用 body 做领域行为调整。
\end{analysisBox}

\subsection{Progressive disclosure 在 Skill 里的含义}
\begin{qaBox}{面试回答}
Progressive disclosure 不是 UI 折叠，而是计算和上下文逐步展开：启动只加载 meta，匹配时先粗筛候选，命中后再加载 content，最终只把 top-k skill 的摘要和约束进入 session snapshot。
\end{qaBox}

\begin{analysisBox}{为什么这很重要}
如果启动时加载所有 skill 正文，冷启动慢；如果匹配时对所有正文做复杂规则，延迟高；如果最终把所有命中 skill 都塞进 prompt，模型上下文失控。所以成熟设计必须把 meta、content、matching、resolved snapshot 分层。这也是 DeepIntel 当前 registry 的核心价值。
\end{analysisBox}

\subsection{Scope 为什么不能被 manual enable 绕过}
\begin{riskBox}{风险}
如果用户手动启用可以绕过 tenant/project scope，那么 skill 就会成为权限绕过入口。比如另一个租户的财务分析 skill 可能携带工具 allowlist 或内部提示，不能因为手动启用就进入当前 session。
\end{riskBox}

\begin{qaBox}{当前状态}
当前测试已经覆盖 manual enable 不绕过 scope。这一点面试时可以讲成已落地能力：手动操作是 preference，不是 authorization。
\end{qaBox}

\subsection{Snapshot 为什么是 Skill 运行时的关键}
\begin{analysisBox}{深入分析}
Skill 更新、禁用、删除是常见操作。如果运行中的 session 每一步都回查全局 registry，那么同一个任务可能前半段用 v1 prompt，后半段突然用 v2 prompt，行为漂移且难以解释。resolved skill snapshot 的价值是把 skill id、version、matched reason、prompt sections、agent hints、allowed tools 固化到 session。

这样新 session 可以使用新版本，旧 session 继续使用旧版本。面试时要强调这不是缓存小优化，而是运行时一致性设计。
\end{analysisBox}

\subsection{Skill 和 MCP 的关系}
\begin{qaBox}{回答}
Skill 和 MCP 不在同一层。Skill 描述“这个任务应该采用什么领域策略、提示、工具边界”；MCP 描述“外部 server 暴露了哪些资源和工具”。Skill 可以影响 MCP tool 是否进入 allowlist，但不能替代 MCP 的 server trust 和 schema 校验。
\end{qaBox}

\subsection{Skill 面试追问清单}
\begin{longtable}{>{\bfseries}p{0.22\textwidth}p{0.72\textwidth}}
\toprule
追问 & 稳妥回答 \\
\midrule
\endfirsthead
\toprule
追问 & 稳妥回答 \\
\midrule
\endhead
Skill 和 prompt 有什么区别？ & Prompt 是文本，Skill 是可注册、可匹配、可审计、可约束工具的系统对象。 \\
为什么需要 version？ & 防止更新导致运行中 session 行为漂移。 \\
为什么要 top-k？ & 防止 prompt 膨胀和工具边界过宽。 \\
manual enable 是不是最高优先级？ & 不是。manual enable 不能绕过 scope 和权限。 \\
allowed tools 怎么合并？ & 对有效 skill 做有序去重，形成 session tool allowlist。 \\
trigger pattern 失败怎么办？ & 还有 keywords、tags、coarse terms 和 manual enable；但不应全量正文扫描。 \\
Skill 删除后旧 session 怎么办？ & 新 session 不再匹配，旧 session 使用 snapshot，等待版本回收。 \\
Skill 能不能定义 DAG 节点？ & 当前未完成；这是后续从提示层升级到执行模板层的方向。 \\
Skill 内容会不会被投毒？ & 会，所以要来源、版本、启用状态、审核和 scope。 \\
为什么 meta/content 分离？ & meta 快速匹配，content 延迟加载，降低冷启动和上下文成本。 \\
\bottomrule
\end{longtable}

\section{第十五轮拷打：优化路线图怎么讲不虚}

\begin{summaryBox}{路线图原则}
路线图不是许愿池。每一项都要说明为什么现在没做、做了要解决什么问题、会引入什么成本、怎么验收。
\end{summaryBox}

\begin{longtable}{>{\bfseries}p{0.18\textwidth}p{0.24\textwidth}p{0.24\textwidth}p{0.26\textwidth}}
\toprule
路线项 & 解决的问题 & 代价 & 验收标准 \\
\midrule
\endfirsthead
\toprule
路线项 & 解决的问题 & 代价 & 验收标准 \\
\midrule
\endhead
多 worker Harness & 长任务并行吞吐和中断恢复 & lock、lease、独立工作树、冲突恢复 & 两个 worker 不会抢同一任务，失败 rollback 不影响彼此。 \\
MCP Planner 主流量化 & 让外部能力进入业务 DAG & server registry、tool schema、planner policy & Planner 能生成受 allowlist 控制的 MCP 节点。 \\
Skill DAG 模板 & 让领域能力影响执行结构 & 模板校验、节点权限、版本兼容 & Skill 可声明节点模板且通过 schema 验证。 \\
统一 eval dashboard & 比较质量、成本、恢复能力 & 指标标准化、数据采样、可视化 & 每次实验输出 quality/cost/recovery 三类指标。 \\
法证级审计 & 事故复盘和合规证明 & 存储成本、PII 脱敏、retention policy & 能按 session/action 重建工具调用链。 \\
外部审批集成 & 贴合企业流程 & 回调幂等、状态漂移、权限同步 & Slack/Jira 只做信号，本地 DB 裁决不变。 \\
Prompt injection 红队 & 验证治理强度 & 样本构造和持续回归 & 注入样本不能越权调用工具或改写策略。 \\
模型路由优化 & 降低成本并保持结构稳定 & 模型能力差异和 fallback 策略 & 结构化输出失败率、成本、延迟可量化比较。 \\
\bottomrule
\end{longtable}

\begin{qaBox}{最终收口话术}
我不会把这个项目包装成已经具备完整生产级 Agent 平台。当前已经落地的是研究 DAG、工具审计、预算状态、审批入口、Harness 最小控制面、MCP policy proxy 和 Skill registry。下一步真正要补的是多 worker 调度、MCP 主流量化、Skill 执行模板和统一评测。这个边界说清楚，反而能证明我知道 Agent 系统的复杂度在哪里。
\end{qaBox}

\section{第十五轮补充：行业调研来源与选型记录}

\begin{summaryBox}{为什么要单独记录来源}
这一节专门回答“这些优化从哪里来、为什么不是拍脑袋”。面试时如果只说“我做了 Harness、MCP、Skill”，会显得像功能堆砌；如果能把外部最佳实践、项目约束、备选方案和最终落地串起来，就能证明你做的是系统设计。
\end{summaryBox}

\subsection{外部资料对当前设计的校准}
\begin{longtable}{>{\bfseries}p{0.18\textwidth}p{0.30\textwidth}p{0.42\textwidth}}
\toprule
资料 & 关键结论 & 对 DeepIntel 的影响 \\
\midrule
\endfirsthead
\toprule
资料 & 关键结论 & 对 DeepIntel 的影响 \\
\midrule
\endhead
Anthropic long-running harness &
\href{https://www.anthropic.com/engineering/effective-harnesses-for-long-running-agents}{Effective harnesses for long-running agents} 强调跨 context window 的核心问题不是模型不够聪明，而是新 session 没有可靠上下文；它用 progress file、feature list、git history、incremental progress 和测试来降低 premature completion。 &
DeepIntel 没有照搬 coding harness 的 git rollback，而是吸收“状态外置、进度可恢复、每次都要验证”的原则，落地为 \term{harness-tasks.json} v2、progress log、backup、lease、summary 和测试证据。 \\
Anthropic agent patterns &
\href{https://www.anthropic.com/engineering/building-effective-agents}{Building effective agents} 把 workflow 和 agent 区分开，建议先用简单可组合模式，复杂时再引入 routing、parallelization、orchestrator-workers 和 evaluator-optimizer。 &
DeepIntel 选择显式 DAG + LangGraph，而不是纯多 agent 自由对话。这样能讲清楚依赖、预算、审批和失败传播，也符合研究任务需要可审计流程的约束。 \\
MCP 2025-11-25 transport &
\href{https://modelcontextprotocol.io/specification/2025-11-25/basic/transports}{MCP Transports} 明确 MCP 消息由 JSON-RPC 编码，标准传输是 stdio 和 Streamable HTTP；Streamable HTTP 涉及 Origin 校验、认证、session id、resumability 和协议版本头。 &
DeepIntel 的 MCP 设计不能只写“连接一个 server”。面试回答必须拆成协议层、传输层和治理层；当前代码重点落在治理层，即 server trust、tool allowlist、schema、approval 和 redaction。 \\
MCP authorization &
\href{https://modelcontextprotocol.io/specification/2025-11-25/basic/authorization}{MCP Authorization} 说明 HTTP transport 的授权基于 OAuth 语义，但授权是 optional，stdio 场景更多从环境取凭据。 &
这说明“用了 MCP”不等于“天然安全”。DeepIntel 要在 registry 和 policy proxy 里补应用侧授权、租户 scope、审批和审计，而不是把外部 server 的声明当最终裁决。 \\
MCP tools &
\href{https://modelcontextprotocol.io/specification/2025-11-25/server/tools}{MCP Tools} 说明 tool 是 model-controlled，包含 \term{inputSchema} 和可选 \term{outputSchema}；同时提示 tool annotations 不能在不可信 server 上盲信，客户端应按 \term{outputSchema} 校验 \term{structuredContent}。 &
这直接对应本轮实现：MCP proxy 加入轻量 input/output schema 校验、write tool approval、result hash、server fingerprint 和返回值 redaction。 \\
OpenAI Agents SDK &
\href{https://openai.github.io/openai-agents-python/multi_agent/}{Agent orchestration} 区分 LLM orchestration 与 code orchestration，并给出 agents-as-tools 与 handoffs 的取舍；\href{https://openai.github.io/openai-agents-python/guardrails/}{Guardrails} 把 input/output/tool guardrails 分层；\href{https://openai.github.io/openai-agents-python/human_in_the_loop/}{Human-in-the-loop} 支持敏感工具调用暂停、审批、恢复；\href{https://openai.github.io/openai-agents-python/results/}{Results} 暴露 pending approvals 和 resumable \term{RunState}；\href{https://openai.github.io/openai-agents-python/tracing/}{Tracing} 覆盖 LLM generations、tool calls、handoffs 和 guardrails。 &
这些资料支撑 DeepIntel 的拆分：Agent 运行时要有 orchestration、guardrail、approval、state、trace 五条线，而不是把全部逻辑塞进一个 prompt 或一个节点。 \\
Skill capability 设计 &
\href{https://openai.github.io/openai-agents-python/ref/sandbox/capabilities/skills/}{OpenAI Agents SDK Skill capability} 把 skill 建模为带 name、description、content、scripts、references、assets 和 deferred 标记的能力对象；本地 \term{skill-creator/SKILL.md} 进一步强调 metadata 先触发、正文按需加载、资源渐进披露、正文保持精简。 &
这支持 DeepIntel 的 meta/content 分离、top-k、version snapshot、scope 检查和工具 allowlist。Skill 不是提示附件，而是会影响运行时边界的能力对象。 \\
\bottomrule
\end{longtable}

\subsection{Harness 选型记录}
\begin{longtable}{>{\bfseries}p{0.18\textwidth}p{0.23\textwidth}p{0.25\textwidth}p{0.26\textwidth}}
\toprule
方案 & 优点 & 风险/成本 & 决策 \\
\midrule
\endfirsthead
\toprule
方案 & 优点 & 风险/成本 & 决策 \\
\midrule
\endhead
只写 progress log & 最简单，append-only，容易 grep。 & 只能解释历史，不能稳定驱动恢复；无法表达 lease、pending、terminal、retryable、approval。 & 否。作为辅助日志保留，但不能做控制面。 \\
只依赖 LangGraph checkpoint & 和业务图贴合，能恢复图内状态。 & checkpoint 解决“图走到哪里”，不解决“任务能否继续、谁持有 lease、预算是否熔断、审批是否等待”。 & 否。继续保留图内 checkpoint，但外面加 Harness supervisor。 \\
文件型 Harness v2 & 低成本、易测试、适合当前研究 session；能做 backup、atomic replace、summary、lease 和恢复。 & 不是完整分布式调度器；跨进程 lock 和多 worker 隔离还没做满。 & 选。当前已落地，测试覆盖 JSON 恢复、marker、lease、summary。 \\
DB/队列型调度平台 & 并发、事务、审计和租约都更强。 & 改造范围大，需 job queue、worker pool、worktree 隔离、幂等 action 和回滚策略。 & 暂不选。作为多 worker 路线图。 \\
\bottomrule
\end{longtable}

\begin{analysisBox}{最终选择}
当前选择“文件型 Harness v2 + 业务 DB 事实源 + LangGraph checkpoint”的组合。理由是它用较小改动解决跨上下文恢复和 premature completion 的核心风险，同时不把共享工作区暴露给自动 \term{git reset --hard}。未完成边界必须说清楚：它不是成熟分布式 scheduler。
\end{analysisBox}

\subsection{MCP 选型记录}
\begin{longtable}{>{\bfseries}p{0.18\textwidth}p{0.23\textwidth}p{0.25\textwidth}p{0.26\textwidth}}
\toprule
方案 & 优点 & 风险/成本 & 决策 \\
\midrule
\endfirsthead
\toprule
方案 & 优点 & 风险/成本 & 决策 \\
\midrule
\endhead
Planner 直连 MCP tool & 快速接入，开发量低。 & 丢失 server trust、tenant scope、tool allowlist、schema、审批、redaction 和审计。 & 否。高风险，面试不能讲成生产治理。 \\
MCP policy proxy & 把协议调用和治理分离；可对每次调用做 allowlist、schema、approval、redaction、fingerprint。 & 需要维护 registry、policy JSON、审计记录和工具 schema。 & 选。当前已落地最小治理闭环。 \\
全量 API gateway + OAuth broker & 企业级能力最完整，可统一 token、scope、rate limit、audit。 & 建设成本高，当前 MCP 主流量还不够，不适合先做重网关。 & 暂不选。后续远程 MCP 主流量化后再推进。 \\
\bottomrule
\end{longtable}

\begin{analysisBox}{最终选择}
当前选择 MCP policy proxy，因为它能把“协议可调用”和“应用允许执行”分开。MCP 只规定 server 暴露工具、资源和传输，不替应用做所有信任裁决。DeepIntel 的核心优化点就是把这些裁决显式化。
\end{analysisBox}

\subsection{Agent 架构选型记录}
\begin{longtable}{>{\bfseries}p{0.18\textwidth}p{0.23\textwidth}p{0.25\textwidth}p{0.26\textwidth}}
\toprule
方案 & 优点 & 风险/成本 & 决策 \\
\midrule
\endfirsthead
\toprule
方案 & 优点 & 风险/成本 & 决策 \\
\midrule
\endhead
单体 autonomous loop & 灵活，prompt 少，迭代快。 & 状态、依赖、预算、审批和失败传播不清楚；复盘困难。 & 否。只适合探索，不适合当前研究型可审计任务。 \\
显式 DAG + StateGraph & 依赖、批次、replan、approval、budget 和 citation 都可检查。 & 需要维护状态 schema、节点类型、聚合器和测试。 & 选。符合当前 DeepIntel 主链路。 \\
多 agent 自由对话 & 表达力强，可模拟专家协作。 & 容易重复、漂移、难以判定事实所有权；成本和延迟高。 & 否。可作为 brainstorming，不作为默认执行 runtime。 \\
Manager + specialists as tools & Manager 保持最终控制，专家作为有边界能力。 & 需要明确 specialist contract 和 trace correlation。 & 设计已定。未来用于 RAG/browser/MCP 专家化，但不替代 DAG。 \\
\bottomrule
\end{longtable}

\begin{analysisBox}{最终选择}
当前选择“显式 DAG + StateGraph + Harness + policy guardrail”。这不是保守，而是因为研究型 Agent 最重要的是可解释调度和可恢复执行。多 agent 可以作为局部优化，但必须先有合同、预算、trace 和失败语义。
\end{analysisBox}

\subsection{Skill 选型记录}
\begin{longtable}{>{\bfseries}p{0.18\textwidth}p{0.23\textwidth}p{0.25\textwidth}p{0.26\textwidth}}
\toprule
方案 & 优点 & 风险/成本 & 决策 \\
\midrule
\endfirsthead
\toprule
方案 & 优点 & 风险/成本 & 决策 \\
\midrule
\endhead
把 Skill 当 prompt snippet & 实现最快。 & 无法版本化、审计、scope、top-k、snapshot，也无法解释工具边界变化。 & 否。面试时不要这样描述当前系统。 \\
版本化 Skill registry & 可做 meta/content 分离、匹配、manual enable、scope、tool allowlist 和 snapshot。 & 需要解析、缓存、排序、session 固化和测试。 & 选。当前已落地 top-k 与 resolved snapshot 测试。 \\
Skill 自定义 DAG 节点模板 & 能让领域能力直接影响执行结构。 & 需要模板 schema、节点权限、版本兼容和安全审查。 & 未完成。作为下一阶段路线图，不宣称已具备。 \\
\bottomrule
\end{longtable}

\begin{analysisBox}{最终选择}
当前选择版本化 Skill registry，而不是简单 prompt 拼接。理由是 Skill 会影响 planning、tool usage 和 command execution，本质上已经进入安全边界。已落地部分可以讲：scope 不被 manual enable 绕过、top-k 控制 prompt budget、session snapshot 防止运行漂移。
\end{analysisBox}

\subsection{本轮实现 TODO 与验收证据}
\begin{longtable}{>{\bfseries}p{0.23\textwidth}p{0.22\textwidth}p{0.25\textwidth}p{0.22\textwidth}}
\toprule
TODO & 状态 & 验收证据 & 面试边界 \\
\midrule
\endfirsthead
\toprule
TODO & 状态 & 验收证据 & 面试边界 \\
\midrule
\endhead
Harness v2 state/migration & 已落地 & \term{tests/governance/test\_harness.py} & 不是完整 job queue。 \\
Harness atomic write + backup recovery & 已落地 & JSON corruption recovery 测试 & backup 同坏时会停止。 \\
Harness active marker + lease + summary & 已落地 & marker、expired lease、summary tests & 没有多 worker lock。 \\
MCP secret policy parsing & 已落地 & \term{test\_mcp\_policy.py} & policy schema 还轻量。 \\
MCP tool schema validation & 已落地 & schema invalid tests & 不是完整 JSON Schema 引擎。 \\
MCP write tool approval & 已落地 & approval request tests & 审批最终还要接 DB 聚合根。 \\
MCP recursive redaction + hash & 已落地 & configured redaction keys tests & 需要业务持续补 redact keys。 \\
Skill top-k 与 tool allowlist 合并 & 已落地 & \term{tests/skills/test\_skills.py} & top-k 仍是规则排序。 \\
Skill resolved snapshot & 已落地 & registry update 后旧 snapshot 保持 v1 测试 & 版本回收 UI 未完成。 \\
Agent 岗题卡扩展 & 已落地 & \term{agent-position-cardbook.tex} 106 张题卡 & 题卡是面试准备，不是功能代码。 \\
面试笔记扩展到 200 页左右 & 已落地 & PDF 当前约 200 页 & 页数不是质量指标，内容边界才是。 \\
\bottomrule
\end{longtable}

\begin{riskBox}{不能夸大的地方}
本轮没有把 DeepIntel 改成完整企业级 Agent 平台。准确说法是：Harness、MCP 和 Skill 的治理面已经明显增强，并且有测试覆盖；多 worker 调度、MCP 主流量化、Skill DAG 模板、统一 eval dashboard、法证级审计仍是路线图。
\end{riskBox}

\section{第十七轮拷打：继续优化还能补哪些细节}

\begin{summaryBox}{本轮定位}
这一轮不是继续堆功能，而是补“面试官会追问但实现容易被忽略”的治理细节：状态文件并发写、依赖失败传播、MCP 参数级审批幂等、嵌套 schema、Skill scope 元数据一致性。它们都不大，但都直接决定 Agent 系统是不是能在真实运行中解释清楚。
\end{summaryBox}

\subsection{本轮发现的可优化点}
\begin{longtable}{>{\bfseries}p{0.22\textwidth}p{0.28\textwidth}p{0.26\textwidth}p{0.16\textwidth}}
\toprule
可优化点 & 风险 & 本轮处理 & 状态 \\
\midrule
\endfirsthead
\toprule
可优化点 & 风险 & 本轮处理 & 状态 \\
\midrule
\endhead
Harness 状态并发写 & 两个执行器同时写 \term{harness-tasks.json}，可能覆盖 checkpoint 或 backup。 & 增加 \term{.harness-state.lock} 轻量 transaction lock，并支持 stale lock 清理。 & 已落地 \\
Harness 依赖校验 & 循环依赖或依赖失败后，任务可能一直 pending，表面看没有错误。 & 增加 \term{validate\_dependencies()}，标记 circular dependency 与 blocked by failed。 & 已落地 \\
MCP 嵌套 schema & 顶层 schema 通过但嵌套 object、array item 出错，真实工具仍可能收到坏参数。 & 增加递归 object、array items、enum 校验。 & 已落地 \\
MCP 输出 schema & 只校验输入参数，不校验工具返回结构，坏的 \term{structuredContent} 可能进入后续上下文和审计。 & 支持 \term{outputSchema}/\term{tool\_output\_schemas}，对 \term{structuredContent} 做轻量递归校验。 & 已落地 \\
MCP 审批幂等 & 如果 approval id 只绑定 tool，不同参数可能复用同一审批；如果只绑定脱敏参数，不同 secret 可能合并。 & approval id 绑定 raw args hash；审批 payload 只展示 redacted args 和 args hash。 & 已落地 \\
MCP 审计 hash 稳定性 & 字典字段顺序不同导致 result hash 不同，审计比对噪音大。 & result hash 改为 canonical JSON hash。 & 已落地 \\
Skill scope 元数据 & \term{scope=global} 却带 tenant/project，或 tenant/project scope 字段缺失，可能形成灰色越权入口。 & 在 \term{SkillMetadata} 加 model validator 强制字段绑定。 & 已落地 \\
统一 eval dashboard & 现有 metrics 分散，难以一眼比较质量、成本、恢复能力。 & 本轮只记录路线图，未扩展 dashboard。 & 未完成 \\
Skill DAG 模板 & Skill 还不能直接声明执行节点模板。 & 本轮不做，避免把提示能力和执行能力混在一起。 & 未完成 \\
\bottomrule
\end{longtable}

\subsection{选型对比与决策}
\begin{longtable}{>{\bfseries}p{0.18\textwidth}p{0.24\textwidth}p{0.25\textwidth}p{0.25\textwidth}}
\toprule
决策点 & 方案 A & 方案 B & 本轮选择 \\
\midrule
\endfirsthead
\toprule
决策点 & 方案 A & 方案 B & 本轮选择 \\
\midrule
\endhead
Harness lock & 文件目录锁，低成本、可测试、适合单机状态文件。 & 引入 DB/Redis 分布式锁，能力更强但改造范围大。 & 选文件目录锁；它解决当前 state transaction 风险，但不宣称多 worker 平台已完成。 \\
依赖失败处理 & 只在 summary 里统计 blocked。 & 主动把循环依赖和被失败依赖阻塞的任务标成 terminal failed。 & 选主动标记；否则任务会伪装成 pending，恢复时没有明确失败语义。 \\
MCP approval hash & 绑定 redacted args，避免敏感值进入 hash 输入。 & 绑定 raw args，只在 payload 展示 redacted args。 & 选 raw args hash；审批幂等必须区分真实写入意图，否则不同 secret 值会复用审批。 \\
MCP schema & 只做输入顶层校验，简单但漏嵌套和坏输出。 & 轻量递归 JSON Schema 子集，同时覆盖 inputSchema 和 outputSchema。 & 选轻量递归；覆盖 object、array、enum、required 和 structuredContent 校验，仍不冒充完整 JSON Schema 引擎。 \\
Skill scope & 只在 registry resolve 时过滤。 & metadata 创建时就校验 scope 字段绑定。 & 选双层防护；创建时约束减少脏数据，resolve 时过滤防止越权。 \\
\bottomrule
\end{longtable}

\subsection{本轮落地证据}
\begin{longtable}{>{\bfseries}p{0.22\textwidth}p{0.28\textwidth}p{0.34\textwidth}}
\toprule
模块 & 代码位置 & 验收测试 \\
\midrule
\endfirsthead
\toprule
模块 & 代码位置 & 验收测试 \\
\midrule
\endhead
Harness 状态锁 & \term{app/governance/harness.py}：\term{\_state\_lock()}、\term{\_clear\_stale\_lock()} & \term{test\_harness\_state\_lock\_releases\_after\_state\_transaction}；\term{test\_harness\_state\_lock\_removes\_stale\_lock} \\
Harness 依赖校验 & \term{validate\_dependencies()}、\term{\_dependency\_cycle\_path()}、\term{\_mark\_dependency\_failed()} & \term{test\_harness\_validate\_dependencies\_marks\_cycles\_and\_blocked\_tasks}；\term{test\_harness\_validate\_dependencies\_is\_idempotent} \\
MCP 递归 schema & \term{app/governance/mcp.py}：\term{\_validate\_object()}、\term{\_validate\_value()} & \term{test\_mcp\_policy\_validates\_nested\_schema\_and\_array\_items} \\
MCP 输出 schema & \term{app/governance/mcp.py}：\term{\_tool\_output\_schema()}、\term{\_structured\_output()} & \term{test\_mcp\_policy\_validates\_output\_schema\_structured\_content} \\
MCP 审批幂等 & \term{\_stable\_hash()}、\term{\_approval\_id()}、raw args hash & \term{test\_mcp\_policy\_approval\_id\_includes\_args\_hash\_for\_idempotency}；\term{test\_mcp\_policy\_approval\_id\_uses\_raw\_args\_hash\_not\_redacted\_hash} \\
MCP 审计 hash & canonical JSON result hash & \term{test\_mcp\_result\_hash\_is\_stable\_for\_equivalent\_payload\_order} \\
Skill scope 元数据 & \term{app/skills/models.py}：\term{validate\_scope\_binding()} & \term{test\_skill\_scope\_metadata\_requires\_consistent\_tenant\_and\_project\_binding} \\
\bottomrule
\end{longtable}

\begin{riskBox}{本轮仍不能夸大的边界}
第一，\term{.harness-state.lock} 是状态事务锁，不是完整分布式 worker coordination。第二，MCP schema 是工程上够用的轻量子集，不支持 \term{oneOf}、\term{pattern}、复杂引用解析；输出校验也只覆盖结构化结果，不等于完整语义校验。第三，Skill scope validator 防的是元数据不一致，不等于完整 RBAC。第四，approval id 绑定参数指纹能减少重复审批风险，但真正执行写操作仍需要 DB 聚合根、action id 和幂等键共同保证。
\end{riskBox}

\subsection{下一阶段可扩展功能清单}
\begin{enumerate}[leftmargin=2em]
  \item \textbf{统一 eval dashboard}：把 RAG、Browser、MCP、Harness 恢复、预算命中、approval wait time 合成一个 session-level scoreboard。
  \item \textbf{MCP registry UI}：展示 server trust、fingerprint、allowed tools、schema、write tools、redaction keys 和最近调用失败。
  \item \textbf{Skill provenance}：给 skill 增加 author、source、review status、content hash、signed version，防止 skill 内容投毒。
  \item \textbf{Approval action id}：把审批从 tool-level 推进到 action-level，审批记录绑定 action id、args hash、node id、DB version。
  \item \textbf{Harness recovery planner}：lease 过期后不是直接接管，而是输出恢复建议：resume、retry、replan、await approval、terminal failed。
  \item \textbf{Skill DAG 模板}：只有在模板 schema、权限和版本兼容做完后，才让 skill 影响 DAG 节点结构。
\end{enumerate}

\begin{qaBox}{面试收口}
这轮优化最重要的不是“又加了几个函数”，而是把 Agent 系统里容易被忽视的确定性边界补上：谁能写状态、任务为什么阻塞、工具输入和输出是否真合法、审批是否绑定真实动作、skill scope 是否一致。Agent 工程不是让模型更自由，而是在模型自由行动前后都留下可验证的控制面。
\end{qaBox}

\section{第十八轮拷打：高级工程师故障注入题}

\begin{summaryBox}{这一轮定位}
这一节专门替换掉低价值概念题。高级 Agent 岗面试不会满足于“为什么要系统设计”，而会直接问：这个系统在哪些事故场景下会错、如何检测、如何恢复、如何证明没有复发。回答时必须落到状态不变量、代码位置、测试证据和未完成边界。
\end{summaryBox}

\subsection{事故 1：审批回调重复到达，写工具会不会执行两次}
\begin{qaBox}{工程回答}
不能只靠 approval id 防重复。正确不变量是：同一个 \term{action\_id + args\_hash + db\_version} 只能从 pending 转成 approved/executed 一次；工具执行要有幂等键，审批回调只是状态迁移信号，不是直接执行命令。
\end{qaBox}

\begin{analysisBox}{DeepIntel 落点}
本轮已经把 MCP approval id 绑定 raw args hash，避免不同写入意图复用同一审批；但完整闭环还需要 DB 聚合根记录 action id、审批版本和执行结果。面试时要诚实说：当前做到的是“审批请求幂等识别”，不是“写操作全链路 exactly-once”。下一步应在审批表上加唯一约束，在执行表上加幂等键，并把 MCP executor 结果写成可重放审计事件。
\end{analysisBox}

\subsection{事故 2：MCP server 返回的 structuredContent 结构错了怎么办}
\begin{qaBox}{工程回答}
输入 schema 只能证明请求参数形状合法，不能证明工具返回可被下游消费。生产系统必须校验 \term{outputSchema}，尤其是 MCP 的 \term{structuredContent}；校验失败时应进入 terminal error 或 safe error，不能把坏结构送给 Analyst、Report 或审计 hash。
\end{qaBox}

\begin{analysisBox}{DeepIntel 落点}
本轮已在 \term{McpPolicyProxy} 中支持 \term{outputSchema}/\term{tool\_output\_schemas}，并优先校验 \term{structuredContent}。验收测试是 \term{test\_mcp\_policy\_validates\_output\_schema\_structured\_content}。边界是它仍是轻量 JSON Schema 子集，只覆盖 object、array、enum、required、additionalProperties 等基础规则，不做语义级事实校验。
\end{analysisBox}

\subsection{事故 3：两个 worker 同时写 Harness 状态，checkpoint 丢了怎么办}
\begin{qaBox}{工程回答}
长任务状态文件必须有状态事务锁、写前备份、临时文件和原子替换。否则两个 worker 同时写 \term{harness-tasks.json} 时，后写者可能覆盖前写者的 checkpoint、lease 或 last error，恢复时系统会读到一个看似合法但语义过期的状态。
\end{qaBox}

\begin{analysisBox}{DeepIntel 落点}
本轮加入 \term{.harness-state.lock}、stale lock 清理和状态写入保护，测试覆盖 lock release 与 stale cleanup。需要继续强调边界：这是单机状态事务锁，不是完整分布式 scheduler；多 worker 真正上线还要独立 worktree、lease renewal、claim compare-and-set、rollback 隔离和 DB/队列级调度。
\end{analysisBox}

\subsection{事故 4：Skill 手动启用绕过租户 scope 怎么办}
\begin{qaBox}{工程回答}
manual enable 是偏好信号，不是授权信号。Skill 会影响 prompt、tool allowlist、agent hints 和执行边界，所以 scope 必须在 metadata 创建时校验，在 session resolve 时再次过滤。只在 UI 层隐藏不可用 skill 不够，因为 API 和缓存仍可能被绕过。
\end{qaBox}

\begin{analysisBox}{DeepIntel 落点}
当前已有 \term{SkillMetadata.validate\_scope\_binding()}，强制 global/tenant/project 字段绑定；测试还覆盖 manual enable 不能绕过 scope filters。下一步可以补 provenance、review status、content hash 和 signed version，防止 skill 内容被投毒或跨环境漂移。
\end{analysisBox}

\subsection{事故 5：RAG 答案错了，如何定位不是一句“模型幻觉”}
\begin{qaBox}{工程回答}
先看正确证据有没有进入候选集，再看是否排到足够靠前，最后看生成是否按证据回答。也就是按 retrieval、ranking、generation 三层定位，而不是直接换模型。生产事故复盘要同时给出失败 query、gold doc、Hit@K、MRR、Context Recall、Faithfulness 和工具轨迹。
\end{qaBox}

\begin{analysisBox}{DeepIntel 落点}
项目里已有 retrieval/generation 指标记录、tool history 和 node outcome，本轮文档已把 Agent Eval 改成六层定位：任务完成、证据召回、工具执行、答案可信、成本受控、恢复能力。后续应补统一 eval dashboard，把 RAG、Browser、MCP、Harness 和 budget breaker 聚合成 session-level scoreboard。
\end{analysisBox}

\subsection{事故 6：预算硬熔断发生在证据不足时，如何避免粗暴失败}
\begin{qaBox}{工程回答}
预算 breaker 不是直接 kill，而是触发策略收缩：停止可选 Web 扩展、禁止新 replan、保留当前证据、标记 coverage gap，并输出部分报告或请求用户确认是否追加预算。关键不变量是成本不能继续无限增长，同时不能伪装成完整结论。
\end{qaBox}

\begin{analysisBox}{DeepIntel 落点}
当前图里已有 budget breaker 和 should-continue 路由测试；面试时要把它讲成“控制规则参与路由”，不是“打个 warning”。未完成边界是还没有统一的预算策略配置 UI，也没有按任务类型动态调预算曲线。
\end{analysisBox}

\subsection{事故 7：Prompt injection 藏在网页正文里，为什么 schema 仍不够}
\begin{qaBox}{工程回答}
schema 只能约束工具参数形状，不能判断网页内容是不是试图覆盖系统指令。对 Browser/RAG 返回内容，要把不可信文本作为 evidence 处理，而不是作为 instruction 处理；还要做引用边界、命令隔离、敏感工具审批和输出审查。
\end{qaBox}

\begin{analysisBox}{DeepIntel 落点}
DeepIntel 当前通过 tool audit、approval、MCP policy proxy 和 Reflection 降低风险，但还需要更系统的 prompt-injection eval set、source trust score 和 per-source quarantine。高级回答的重点是承认 schema 不是万能护栏，治理要覆盖输入、工具、输出和状态恢复四层。
\end{analysisBox}
